\documentclass[10pt,a4paper]{amsart}
\usepackage[margin=19mm]{geometry}
\usepackage{amsmath,amssymb,mathtools}
\usepackage[scr=rsfs,cal=euler]{mathalfa}
\usepackage{xfrac}
\usepackage[unicode,hidelinks,bookmarks=false]{hyperref}
\newcommand{\ie}{i.e.\ }
\newcommand{\cf}{cf.\ }
\newcommand{\resp}{respectively\ }
\newcommand{\Subsection}{\subsection}
\providecommand{\nobreakdash}{\nobreak\hbox{-}}
\setlength{\emergencystretch}{2em}
\multlinegap0pt
\usepackage{bm}
\usepackage{bbm}
\usepackage{dsfont}
\usepackage{xspace}
\usepackage{cancel}
\usepackage{mathtools}
\usepackage{amsthm}
\makeatletter
\@ifundefined{equation}{\newtheorem{equation}{Equation}}{}
\@ifundefined{prop}{\newtheorem{prop}[equation]{Proposition}}{}
\@ifundefined{thm}{\newtheorem{thm}[equation]{Theorem}}{}
\makeatother
\makeatletter
\newcommand{\GF}{\mathrm{GF}}
\newcommand{\Map}{\mathrm{Map}}
\newcommand{\R}{\mathbf{R}}
\expandafter\ifx\csname SC\endcsname\relax
\newenvironment{SC}{\color{darkgreen}{\bf VS:} \footnotesize}{}
\fi
\expandafter\ifx\csname VS\endcsname\relax
\newenvironment{VS}{\color{blue}{\bf VS:} \footnotesize}{}
\fi
\newcommand{\del}{\partial}
\makeatother
\title[A topological classification of generating functions]{A topological classification of generating functions}
\author{Sylvain Courte, Vivek Shende}
\date{Source: arXiv:2606.00687v1 (CC BY 4.0)}
\begin{document}
\maketitle

\noindent\emph{Source and licence.} A topological classification of generating functions. Version arXiv:2606.00687v1 (CC BY 4.0), distributed under the Creative Commons Attribution 4.0 International licence, as recorded in the arXiv licence metadata for this exact version and shown on the abstract page https://arxiv.org/abs/2606.00687v1. Full rights notice: licenses/arxiv-2606.00687-CC-BY-4.0.txt.\par\smallskip

\noindent\textit{Editorial scope.} Counted unit: the Thm whose source label is \texttt{thm:GFdoubling} in the article (source0.tex lines 1024--1026, source label \texttt{thm:GFdoubling}), delivered together with the article's own complete original proof of it (source0.tex lines 1027--1058, one \texttt{proof} environment of 32 lines). No mathematics is written, repaired or re-derived: statement and proof are the article's own text line for line. Required context retained uncounted: prop:emptyleg (lines 635--637). Every definition, display and label of those lines is retained unchanged. Omitted from this package and not redistributed: 1--634, 638--1023, 1059--2043. Nothing inside a retained excerpt is omitted or altered: every retained body is delivered exactly as the article's lines are written, so no omission receipt of any kind accompanies this package. Citations: the retained excerpts carry no citation command at all, so no bibliography entry of the article is transcribed and no reference is redistributed. Reference adaptation: none was needed and none was made; every reference inside the delivered unit resolves inside it, so no unresolved reference or citation is produced. Printed slips kept literal: no printed word, symbol, hypothesis or number of the counted statement or of its proof was altered. Layout and class: the article's own class is \texttt{article} with packages including fullpage, inputenc, fontenc, amsmath, amsfonts, amssymb, graphicx, amsthm, xy, tikz-cd, say, hyperref, tocloft; the wrapper is the lane's pinned \texttt{amsart} editorial wrapper at 10pt on A4, which restates the theorem-like environments the retained text needs and adds the article's own macro definitions that the retained text needs (\texttt{\textbackslash GF}, \texttt{\textbackslash Map}, \texttt{\textbackslash R}, \texttt{\textbackslash del}), with extra wrapper packages (bm, bbm, dsfont, xspace, cancel, mathtools). The delivered numbering is the unit's own. Assets: no retained span contains a figure environment, a graphics-inclusion command, a PDF-page inclusion, a table or a TikZ picture; no image file is copied into this package, the package declares no asset dependency and no image file is redistributed with it. Licence: version arXiv:2606.00687v1 of this article is distributed under the Creative Commons Attribution 4.0 International licence, as recorded in the arXiv licence metadata for this exact version and shown on the abstract page https://arxiv.org/abs/2606.00687v1. A full rights notice is delivered at licenses/arxiv-2606.00687-CC-BY-4.0.txt. Characters: every retained excerpt is pure ASCII, so the delivered engine, XeLaTeX with its default Latin Modern text font, reports no missing character.
\begin{prop}\label{prop:emptyleg}
Let $M$ be a manifold and $\Lambda$ the empty Legendrian submanifold of $J^1(M)$. Then $\GF_\Lambda(M)^*$ is contractible.
\end{prop}

\begin{thm}\label{thm:GFdoubling}
The map $\GF^{\leq 0}_{\Lambda\times \Lambda_D}(M\times \R)^* \to \mu\GF^{\leq 0}_{\Lambda\times \Lambda_D}(M\times \R)^*$ is a homotopy equivalence.
\end{thm}

\begin{proof}
We first consider the case where $\Lambda$ is empty. In this case, the right hand side is reduced to a single point so we need to prove
that the left hand side is contractible.
Since $\Lambda$ is empty, the restriction map at $t=0$ (or at any $t=t_0$ for $t_0\leq 0$) $\GF^{\leq 0}_\emptyset(M\times \R) \to \GF_\emptyset(M)$
is a homotopy equivalence. The result now follows from Proposition~\ref{prop:emptyleg}.

Another extreme case is when $M$ is a single point and $\Lambda$ consists of finitely many points $z_1,\dots,z_k$ in $\R=J^1(*)$.
Let $\phi\in \mu\GF^{\leq 0}_{\Lambda\times \Lambda_D}(\R)^*$, this is a germ of function $\phi:\R^n\times \R \to \R$
near finitely many points $(v_1,0),\dots,(v_k,0)$. For each of these points, we pick an embedding $\psi_i:D^{n-1} \times [z_i-a,z_i+a]\times[-\epsilon,\epsilon]\to \R^n\times \R$ whose image is contained in the domain of definition of $\phi$ ($a>0$ and $\epsilon>0$ can be taken arbitrarily small) and satisfies:
$\psi_i(0,z_i,0)=(v_i,0)$, $\psi_i(x,z,t)=(j(x,z,t),t)$, $\phi(\psi_i(x,z,-\epsilon))=z$, $\phi(\psi_i(x,z_i\pm a,t))=z_i\pm a$
and $\phi(\psi_i(x,z,t))=z$ for $|x|=1$. Denote $K_i$ the image of $\psi_i$.
We need to prove that the space of functions $f:\R^n\to \R$ which coincide with $\phi$ on each $K_i$ is highly connected (thus contractible in the limit $n\to \infty$).
We claim that restriction at $t=-\epsilon$ yields a homotopy equivalence between this space and the space of functions
$f:\R^n\to \R$ which have no critical points and agree with $\phi$ on $K_i\cap\{t=-\epsilon\}$, namely $f(\psi(x,z,-\epsilon))=z$.
Indeed by a suitable retraction of $\{t\leq \epsilon\}\subset \R^n\times \R$ onto $\{t\leq -\epsilon\}\cup \cup_i K_i$, we can deform
such functions until they coincide on $\{t\leq \epsilon\}$ provided they already coincide in $\{t\leq-\epsilon\}\cup \cup_i K_i$.
Finally we need to argue that the space $F$ of functions $\R^n\to \R$ without critical points which coincide with the projection $D^{n-1}\times [z_i-a,z_i+a]\to [z_i-a,z_i+a]$ in finitely many such disjoint cylinders embedded in $\R^n$, is highly connected when $n$ goes to infinity.
We claim first that this space is equivalent to the space of functions without critical points $f$ such that $\nabla f(v_i)=dz$.
Indeed, if $\nabla f (v_i)=dz$, we can deform slightly $f$ so that $f=z$ near $v_i$ (the center of $K_i$), then by dilating $K_i$ we can arrange $f=z$
on the whole $K_i$. From the space $E$ of all functions $\R^n\to \R$ without critical points, the collection of $\nabla f(v_i)$ thus defines a map to $(\R^n\setminus\{0\})^k\simeq (S^{n-1})^k$ whose fiber is homotopy equivalent to the space $F$.
The space $(S^{n-1})^k$ is $(n-2)$-connected and $E$ is also highly connected (see Proposition~\ref{prop:emptyleg}), therefore $F$ is highly connected.

We now turn to the general case. We consider the fiber over some semi-local generating function $\phi$ for $\Lambda\times \Lambda_D$,
namely the space of generating functions $f:M\times \R\times \R^n\to \R$ for $\Lambda\times \Lambda_D$ over $M\times (-\infty, \epsilon)$
which agree with $\phi$ near $\Sigma_\phi=\Sigma_f$. Denote $\Sigma_0=\Sigma_\phi\cap\{t=0\}$ and $\phi_0$ the restriction of $\phi$ to $\Sigma_0$.
We pick a tubular neighborhood $U\times [-a,a]$ of $\Sigma_0$ (denote $r:U\to \Sigma_0$ the corresponding projection) and
an embedding $\psi:U\times[-a,a]\times[-\epsilon,\epsilon]\to M\times \R^n\times \R$ such that
$\psi$ is the identity on $\Sigma_0\times\{0\}$, $\psi(u,z,t)=(*,*,t)$, $f(\psi(u,z,-\epsilon))=\phi_0(r(u))+z$ and $f(\psi(u,z,t))=\phi_0(r(u))+z$ for $u\in \del U$.
As in the previous case, the function $f$ is then determined up to contractible choice by its restriction to $t=-\epsilon$. The gradient $\nabla f$
along $\Sigma_0$ defines a map from the space of functions $f:M\times \R^n\to \R$ fiberwise without critical points to
the space $\Map(\Sigma_0, S^{n-1})$, both of which are highly connected, with fiber homotopy equivalent to the space we are considering, so the claim follows.
\end{proof}

\end{document}
